\section{Kiss me, my mortal lover}

In which are discussed encounters with demons, animal spirits, the guardian of the threshold, and dreams of the Anima.

\paragraph{Connecting with demons}

In a recent comment, ``Tom" writes:

\begin{quotex}
At various points in the Inferno, Dante, relying on Virgil, is able to enlist the help of numerous demons, such as the
Centaurs and Titans. Enlisting their help is a terrifying experience for Dante, but with Virgil's
ability to command (or at least negotiate with) them, the threat they pose is neutralized, and they are instead shown
to be indispensable auxiliaries in the journey through the underworld.

The interpretation I made, based on some of my own experiences, is that Dante is describing how in our earthly life, we
often need to make use of these ``lower" powers while remaining firmly in command of them. An example that comes to mind
is in physical battle — a man would find it difficult to fight without awakening the fierce
Centaur nature, the key being then to remain firmly in command of it. 

\end{quotex}
That is why the Templars were accused of devil worship; you cannot be a warrior-monk without the power of these lower
forces. The Templars were an initiatic group, so outsiders easily misunderstand their aims and methods. As Valentin
Tomberg pointed out, the Our Father prayer penetrates from above to the heart\footnote{\url{https://www.meditationsonthetarot.com/protection-meditation}}. Thus it affects the intellect and the
loving-kindness of the heart. That is fine for a monk, but insufficient for a warrior. At its worse that attitude leads
to intellectualism and sentimentalism. A warrior needs the prayer to penetrate down through the lower centres, which
also need to be redeemed. But without the higher centres, the warrior is just a brute. The solution is self-mastery, as
Tomberg makes clear:

\begin{quotex}
One need not fear the devil, but rather the perverse tendencies in oneself! For these perverse human tendencies can
deprive us of our freedom and enslave us. 

\end{quotex}
Otherwise, these perverse tendencies are projected so that they appear to have an ``I" opposed to our own conscious ego.
Jung explains how that works psychologically:

\begin{quotex}
It appears as an autonomous formation intruding upon consciousness… It is just as if the complex were an autonomous
being capable of interfering with the intentions of the ego. \flright{\emph{Psychology and Religion}}

\end{quotex}
\paragraph{Adam and the Animals}
In Eden, there is no reason to believe that animals did not act differently from how we know them now. This is also the
position of Thomas Aquinas; since the animals have no possibility of participating in the supernatural, they could not,
ipso facto, have the gift of eternal life. Besides, animals have a certain essence which is exemplified in their
anatomy and physiology.

In this context, the idea came up that perhaps in Eden, Adam existed at a higher level than the physical. It is true
that man knows animals in essence, not just sensually. To hold the essence of an animal in consciousness, is tantamount
to being that animal. Obviously, this does not mean sensually as in shape-shifting, but rather psychologically. In the
Philokalia, and other spiritual writings, there are may allusions to animals representing various interior states. St
Basil understands the command to ``rule" the animals as the self-mastery of the Intellect over these animalistic states.

These are the states mentioned above, as well as in the Tomberg link. The lower centres are experienced as fear (root
chakra), sex (sacral chakra), and hunger (solar plexus). These are the primary motivations for animal life. These are
various manifestation of the libido, as described in \emph{Symbols of Transformation}.

\paragraph{Dangers of Meditation}
Meditation has reached into pop culture, primarily as a means to achieve beneficial social and life effects, under names
such as mindfulness and so on. Even corporate America is promoting mindfulness in the hope of creating better
employees. However, more recently it is becoming apparent that it is not all peaches and cream. Often times, these
meditators having negative experiences as suppressed parts of the psyche suddenly come into consciousness.

These meditations are conducted without proper instruction or supervision. Worse, what is called ``meditation" is
actually concentration, the first stage that precedes actual meditation.

\paragraph{The Test}
Recently, we received this plaint: I have some trouble with the weekly task of knowing myself. …
doesn't it get lonely? A neophyte often reaches this point; he can choose to move forward or he
can back down. Actually, these are relatively simple tasks. In the novel \emph{Zanoni}, Mejnour prepares Glyndon for
initiation, who then fails the tasks assigned to him:

\begin{itemize}
\item Glyndon gets involved with a local maiden, exciting his sensual appetites 
\item He ignores a beggar in need of alms, demonstrating his stinginess 
\item He decides to go his own way, ignoring the command to not read the book 
\end{itemize}
Glyndon's mind goes in circles until he is finally able to justify his choices. He forgot that the
task is to purify the soul, not to accumulate enjoyable experiences. So, like Glyndon, the lonely fellow gave it all
up, from where he can neither advance nor retreat.

\paragraph{Dweller of the Threshold}
Just like the corporate meditators, the unprepared neophyte will have to confront the lower forces. The novel
\emph{Zanoni} describes Glyndon's experience:

\begin{quotex}
The cloud retreated from it as it advanced; the bright lamps grew wan, and flickered restlessly as at the breath of its
presence. Its form was veiled as the face, but the outline was that of a female; yet it moved not as move even the
ghosts that simulate the living. It seemed rather to crawl as some vast misshapen reptile; and pausing, at length it
cowered beside the table which held the mystic volume, and again fixed its eyes through the filmy veil on the rash
invoker. All fancies, the most grotesque, of monk or painter in the early North, would have failed to give to the
visage of imp or fiend that aspect of deadly malignity which spoke to the shuddering nature in those eyes alone. All
else so dark—shrouded, veiled and larva-like. But that burning glare so intense, so livid, yet so
living, had in it something that was almost HUMAN in its passion of hate and mockery.

The Image spoke to him: his soul rather than his ear comprehended the words it said.

``Thou hast entered the immeasurable region. I am the Dweller of the Threshold. What wouldst thou with me? Silent? Dost
thou fear me? Am I not thy beloved? Is it not for me that thou hast rendered up the delights of thy race? Wouldst thou
be wise? Mine is the wisdom of the countless ages. Kiss me, my mortal lover." 

\end{quotex}
\paragraph{Dreaming the Anima}
For Glyndon, the guardian was a terrifying female figure, who ultimately led him into unconsciousness. That is why
dreams of the Anima need to be recorded and remembered. Here is a test case I know about, which took place over more
several decades. I'll try to use his own words, the best I can recall them.

\subparagraph{Lady of the Lake}
Although I had often had dreams that felt so real, I woke up convinced the woman was real. In the twilight state of
awakening, I would scan my mind to recall where and when I had met her. However, my first awareness of the Anima was a
recurring dream about a beautiful woman, lying face up, unconscious, in a shallow stream. Researching this, I came to
realize that the image is not so uncommon. The myth of the Lady of the Lake had taken many forms, and there have been
paintings of scenes much like the one in my dream.

It seems clear now that the water represents the unconscious and the Anima is waiting to become a conscious part of
myself.

\subparagraph{The Reality Show}
A later dream revealed the purpose of the Anima in my psyche. I was in the audience and on the stage was the emcee. He
announced that he was running a ``reality show" and asked for volunteers from the audience. A woman came up to me and
led me to the stage. I understood the incident to mean that the Anima would lead me to reality.

I commented to him that often the dreamer, who believes he is a spectator in his own dream, turns out to be the ``star"
of his dream instead.

\subparagraph{Lady in attic}
Years later, he reported a different type of dream. Although there were others, this one stands out. In his own words: I
found myself in an old, multi-storied, dark house with narrow stairways. I worked my way up, until I got to the top
floor. A woman then took hold of me and led me up a ladder to a loft. There she opened an old chest, full of antique
books. Excited, I try to read and remember as much as I could. But then I woke up, having forgotten it all.

Yet, it was clear to me that he hadn't forgotten it all. His life changed as he went in search of
antique wisdom. And perhaps it really happened as described.

\subparagraph{The Sordid Affair}
The final dream happened just this week. It was so unlike his other dreams, that he called me immediately the next
morning in a state of confusion.

\subparagraph{The Reunion}
\textit{UPDATE 16 October 2021}

The most recent dream, from a few months ago, involved a union with the Anima. She was under attack by two shady
characters, a man and a woman, and I found myself in the position of protecting her from their attempts to kidnap her.
This aligned with my personal life at the time.



\flrightit{Posted on 2020-10-16 by Cologero }

\begin{center}* * *\end{center}

\begin{footnotesize}\begin{sffamily}



\texttt{Logres on 2020-10-19 at 19:53 said: }

Someone, maybe Steiner, made the point that the Templars were tortured, which released a projection of their psyche.
This should have been more to their credit, as we don't hear that happening much today (and
torture hasn't stopped): at least they were capable of crystallizing an image of evil against
which they were struggling, bringing it to bay. One could make an analogous argument about dreams, and the ability to
crystallize something in dreams, which can then be examined later (if recalled) in light of daytime consciousness
(bringing day and night together, through memory). Thank you for the post. Someone else made a comment about the Grail,
since you bring up the Lady of the Lake, that it was the union of the solar, telluric currents birthing the lunar
current. Mouravieff talks about the moon being a ``fetal earth". Our head is astrologically a ``moon". The solar and
earth currents, then, in the heart and from the lower chakras, have to be brought up into the head (their energies
anyway), and made not just conscious, but super-conscious. It's interesting that when people think
of calm water, they often think of the moonlight on the water. At least I do. Perhaps one could say poetically that
memory is the lunar current, and has to become Reality. It unites night and day, if strong. And Plato talks a lot about
``remembering". So there is that as well.


\hfill

\texttt{Sensitive Ears on 2020-10-22 at 23:09 said: }

That story's obscene, I advise a NSFW warning. IMO you shouldn't share such
dreams. Aren't spiritual success stories more inspiring?


\hfill

\texttt{Cologero on 2020-10-24 at 14:23 said: }

Well, Sensitive, please don't be like that woman on Amazon who reviewed a
children's Bible. She was aghast because God got so angry at Adam and Eve within the first few
pages! Nevertheless, you bring up a good point that I've often pondered. In our age, we have
become accustomed to vulgar language. On the one hand, that is unfortunate. On the other hand, though, older
bowdlerized spiritual literature may be misunderstood. This project is aimed in particular to those who can neither
relate to the modern world nor to the often saccharin and moralistic way that spiritual life is often depicted. It is
not at all the way Joel Osteen presents it. Rather, it is a struggle, an inner combat, a process of purification. Is it
helpful to describe such purifications in explicit detail? I think perhaps so. That's because many
might think that such temptations are unusual or that they represent failure; they may be ashamed to admit to them. The
hermeneutics of suspicion teach that such depictions arise from the ``real" you. We reject that notion. The ``real" you
is something much higher and transcendent.

Keep in mind that the dream in question was quite disturbing to my friend. Yet he should learn an important lesson from
it. For now, I'll follow your suggestion and put similar content behind a password protected page.


\hfill

\texttt{Sensitive Ears on 2020-10-25 at 23:30 said: }

Thanks for the NSFW wall. Some young readers might automatically put themselves in the
protagonist's body and accidentally commit fornication of the heart.


\hfill

\texttt{Mike on 2020-11-10 at 04:49 said: }

To Miss/Mr. Sensitive, I hope you never come across the more ``vulgar" works of Boccaccio, nor the unapologetic, raw
depictions in the different Tantric lineages, whose intensity sieve out the cowardly, the prudish and the base who are
easily distracted and fail to see the reflection of the spiritual bliss inherent in daily pleasure. Your consciousness
rushes to points of pleasure, whether it's induced by tasty ice-cream, or the the ecstasy in
sexual union. All other thoughts are abandoned in such pleasure and you are in a way, reduced to a point of
concentration through these natural processes. For those who know how to deepen that reduction, then, they go beyond
their carnal limits and the ``base act" is at once transformed into a participation in the ``beatific vision"…as above,
so below, only upside down…that's all I'll say as a little hint.

Old man Salvo will probably censor this, or give some snark remark to it lmao. But like the election results, maybe he
will shock me!


\hfill

\texttt{Jupiter on 2022-10-16 at 19:21 said: }

Love and look after your cat that's what they do! Protect ones sleeping aether world and do no harm


\end{sffamily}\end{footnotesize}
